\subsection*{21 października 2026 - przygotowanie projektu}
\textbf{Narzędzie i model:} Codex, rodzina GPT-6; dokładny wariant modelu
nie został udostępniony w interfejsie tej sesji. Korzystano także z podagentów
w ramach Codex do analizy wymagań i przygotowania narzędzia dziennika.

\textbf{Cel:} odczytanie dwóch dokumentów kursowych i przygotowanie
repozytorium oraz dokumentacji wymaganej do dalszej realizacji projektu.

\textbf{Zakres pomocy:} analiza wymagań P00 i P01, zestawienie terminów
i list kontrolnych, przygotowanie struktury repozytorium i szablonu raportu
LaTeX, narzędzia zapisu promptów i odpowiedzi oraz sprawdzenie dokumentacji.
Na tym etapie AI nie opracowało algorytmu RL, środowiska ani wyników eksperymentów.

\textbf{Weryfikacja studenta:} niepotwierdzona. Testy i kontrole wykonane
przez agenta nie stanowią oświadczenia o samodzielnej weryfikacji studenta.
Po zapoznaniu się z materiałami student powinien wskazać, które elementy
sam sprawdził, zrozumiał lub zmodyfikował.

\textbf{Transkrypt:} pełne widoczne prompty i odpowiedzi są zapisywane lokalnie
w dzienniku \texttt{private/ai-journal/main.tex} wraz z tekstami źródłowymi
i metadanymi. Sposób udostępnienia transkryptów prowadzącemu wymaga ustalenia.
Ta rozmowa odbywa się w narzędziu programistycznym; wymaganie dodatkowego
linku dla rozmów poza takim narzędziem nie ma tu zastosowania.

\textbf{Dalsze aktualizacje:} każdą kolejną rozmowę dotyczącą projektu należy
zapisać w całości oraz uzupełnić niniejszą sekcję o datę, model, cel, zakres
wykorzystania i rzeczywisty status samodzielnej weryfikacji studenta.
